\documentclass{article}
\usepackage{graphicx} %for images
\usepackage{amsmath}
\usepackage{amssymb}
\usepackage{listings}
\usepackage[margin=1in]{geometry}
\usepackage[colorlinks=true, allcolors=blue]{hyperref}

\title{Modeling Travel in Gravitational FieldsModeling Travel in Gravitational Fields}
\author{Brian Zhou, Minh Nguyen, James Overcash} 
\date{April 2026}

\begin{document}

\maketitle

\section{Introduction}
A rocket with some known initial position $p_0$, velocity $v_0$, and mass $M$ is traveling through space heading in the direction of $v_0$. The object is propelled by thrusters with thrust $f_t$ and exhaust speed $v_t$. In space lie $N$ planets with fixed positions $p_1, p_2, \dots, p_N$ and known masses $m_1, m_2, \dots, m_N$  whose gravity influences the object's motion. Can we determine whether the rocket will ever land on a target planet?

\section{Analysis}
We model the scenario using classical Newtonian mechanics. We consider the rocket and planets to all be point-masses. Thus, according to Newton's second law, the acceleration $a$ of the object can be related to the net force exerted on the rocket $F_{\text{net}}$ by the equation
\[
\vec{F_{\text{net}}} = m \vec{a}.
\]
We will also assume that space is a perfect vacuum. Thus, the rocket only experiences a gravitational force due to the planets and a thrust force due to it's thruster. According to Newton's law of universal gravitation (Neil Gehrels Swift Learning Center, n.d.), the $i$th gravitational force $F_i$ on the rocket due to the $i$th planet is of the form
\[
\vec{F_i}(t) = G\frac{M m_i}{\lvert \vec{p(t)} - \vec{p_i} \rvert^3}(\vec{p(t)} - \vec{p_i}),
\]
where $G \approx 6.674 \times 10^{-11}$ is Newton's gravitational constant.

The non-linear system is that of the gravitational acceleration. $\vec{F}$ is the force on the object by a mass $m_1$ where the displacement between them is $\vec{r}$. 
\begin{align*}
          F &= G\frac{m \cdot m_1}{ |\vec{r}| ^ 2} \cdot \hat{r} \\ 
    \vec{a} &= G_x\frac{m_1}{|\vec{r}|^3} \cdot \vec{r} \\
    \begin{pmatrix} x'' \\ y'' \end{pmatrix} &= \begin{pmatrix}G\frac{}{}\end{pmatrix}
\end{align*}




\section{Conclusion}
We use Euler's method to quickly simulate paths using the Python code attached in the appendix. 
We find that some limitations of the computational approach includes: we cannot guarantee that the rocket will not eventually land on a planet, we cannot guarantee the path found is the optimal path.
\section{References}
Result in Figure 1 were obtained using a Python script provided in Appendix 1
\begin{thebibliography}{99}
\bibitem{golbeck2015}
Neil Gehrels Swift Learning Center. (n.d.). Imagine.gsfc.nasa.gov. \url{https://imagine.gsfc.nasa.gov/observatories/learning/swift/classroom/law_grav_derivation.html}
\end{thebibliography}


\section{Appendices}
`model.py` contains the code modeling the trajectory of the rocket:
\begin{lstlisting}[language=Python]
from interfaces import *
import tqdm
from os.path import join
import datetime

class Model:
    def __init__(self, rocket: Rocket, 
                 planets: Sequence[Planet], 
                 simulation_seconds: int = 1000 * 3600, 
                 dt: int = 10):
        self.rocket = rocket
        self.planets = planets
        self.path = []
        self.simulation_seconds = simulation_seconds
        self.dt = dt  # units: seconds

    def save_path(self, path: str | None = None):
        SAVE_PATH = "paths"

        if not self.path:
            raise ValueError("Path not created")

        now = datetime.datetime.now()
        if path is None:
            path = join(SAVE_PATH, f"{now.month}_{now.day}_{now.hour}_{now.minute}_{now.second}.path")

        with open(path, mode="w") as f:
            f.write(str(self.dt) + "\n")
            f.write(str(len(self.planets)) + "\n")
            for planet in self.planets:
                f.write(str(planet) + "\n")
            for position in self.path:
                f.write(str(position) + "\n")

    def estimate_path(self):
        for _ in tqdm.tqdm(range(self.simulation_seconds//self.dt)):
            self.rocket.velocity += self.calc_force() / self.rocket.mass_ship * self.dt
            self.rocket.position += self.rocket.velocity * self.dt
            self.path.append(self.rocket.position)

    def estimate_path_with_mass_update(self):
        for _ in tqdm.tqdm(range(self.simulation_seconds//self.dt)):
            thrust = Vector2d(0,0)
            if self.rocket.mass_fuel > 0:
                thrust = self.rocket.thrust
                """
                dM * speed_fuel - M dV = 0
                since thurst = | M dV/dt |
                dM * speed_fuel = | thurst | * dt
                """
                dmassfuel = - thrust.mag() / self.rocket.speed_fuel * self.dt
                self.rocket.mass_fuel += dmassfuel

            self.rocket.velocity += (self.calc_force() + self.rocket.thrust) / (self.rocket.mass_ship + self.rocket.mass_fuel) * self.dt
            self.rocket.position += self.rocket.velocity * self.dt
            self.path.append(self.rocket.position)

    def calc_force(self) -> Vector2d:
        force = Vector2d([0, 0])
        for planet in self.planets:
            dist = max(self.rocket.position.dist(planet.position), 1e-3)
            magnitude = G * self.rocket.mass_ship * planet.mass / dist**3
            force += (planet.position - self.rocket.position) * magnitude
        return force


if __name__ == "__main__":
    r = Rocket(Vector2d(0, 0), 100, Vector2d(0, sqrt(10 * G)))
    planets = [
        Planet(Vector2d(10, 10), 100, 10),
        Planet(Vector2d(0, 10), 100, 10),
        Planet(Vector2d(10, 0), 100, 10),
    ]
    m = Model(r, planets)
    m.estimate_path()
    m.save_path()
\end{lstlisting}

`interfaces.py` contain the interfaces used in `model.py`:
\begin{lstlisting}
from dataclasses import dataclass
from typing import overload
from collections.abc import Sequence, Callable
from math import sqrt


class Vector2d:
    pos: tuple[float, float]

    @overload
    def __init__(self, x: float, y: float, /) -> None: ...

    @overload
    def __init__(self, pos: Sequence[float,], /) -> None: ...

    def __init__(self, *args) -> None:
        if len(args) > 2:
            raise ValueError()
        if len(args) == 1:
            self.pos = tuple(args[0])
        elif len(args) == 2:
            self.pos = args

    def dist(self, other: "Vector2d"):
        return sqrt((self.x - other.x) ** 2 + (self.y - other.y) ** 2)

    def mag(self):
        return sqrt(self.x ** 2 + self.y ** 2)

    @property
    def x(self):
        return self.pos[0]
    
    @x.setter
    def x(self, val: float):
        self.pos = val, self.pos[1]

    @property
    def y(self):
        return self.pos[1]

    @y.setter
    def y(self, val: float):
        self.pos = self.pos[0], val

    def __str__(self) -> str:
        return f"{self.x} {self.y}"

    def __add__(self, other: "Vector2d"):
        return Vector2d(self.x + other.x, self.y + other.y)

    def __sub__(self, other: "Vector2d"):
        return Vector2d(self.x - other.x, self.y - other.y)

    def __mul__(self, other: float):
        return Vector2d(self.x * other, self.y * other)

    def __rmul__(self, other: float):
        return self * other

    def __truediv__(self, other: float):
        return Vector2d(self.x / other, self.y / other)


@dataclass
class Rocket:
    position: Vector2d
    mass_ship: float
    velocity: Vector2d

    mass_fuel: float = 0
    speed_fuel: float = 0
    thrust: Vector2d = Vector2d(0, 0)


@dataclass
class Planet:
    position: Vector2d
    mass: float
    radius: float = 0
\end{lstlisting}


\section{Websites}

\end{document}